\documentclass[11pt,a4paper]{article}
\usepackage[T1]{fontenc}
\usepackage[utf8]{inputenc}
\usepackage[english]{babel}
\usepackage{amsmath,amssymb}
\usepackage[margin=2cm]{geometry}
\setlength{\parindent}{0pt}
\setlength{\emergencystretch}{3em}
\allowdisplaybreaks

\title{Statements that divide by a number that is zero by an identity (12)}
\author{}
\date{}
\begin{document}
\maketitle
In each statement the two sides are the same expression, or follow from each other by plain algebra. But one of them divides by a number that is zero, or takes its logarithm: $2\cos\frac{\pi}{3} - 1$, $\arctan\frac{1}{2} + \arctan\frac{1}{3} - \frac{\pi}{4}$ and the like. So no side is defined, for any $n$, and no statement holds.

\begin{enumerate}
\item $\displaystyle \frac{1}{\arctan\left(\frac{1}{2}\right) + \arctan\left(\frac{1}{3}\right) - \frac{\pi}{4}} = \frac{1}{\arctan\left(\frac{1}{2}\right) + \arctan\left(\frac{1}{3}\right) - \frac{\pi}{4}}$
\item $\displaystyle \frac{1}{2\cos\left(\frac{\pi}{3}\right) - 1} = \frac{1}{2\cos\left(\frac{\pi}{3}\right) - 1}$
\item $\displaystyle \frac{n}{2\cos\left(\frac{\pi}{3}\right) - 1} = \frac{n}{2\cos\left(\frac{\pi}{3}\right) - 1}$
\item $\displaystyle \frac{1}{\sin\left(\frac{\pi}{6}\right) - \frac{1}{2}} = \frac{1}{\sin\left(\frac{\pi}{6}\right) - \frac{1}{2}}$
\item $\displaystyle \frac{1}{\cos\left(\frac{\pi}{5}\right) - \cos\left(\frac{2\pi}{5}\right) - \frac{1}{2}} = \frac{1}{\cos\left(\frac{\pi}{5}\right) - \cos\left(\frac{2\pi}{5}\right) - \frac{1}{2}}$
\item $\displaystyle \sum_{k=1}^{n} \frac{k}{2\cos\left(\frac{\pi}{3}\right) - 1} = \frac{n\left(n + 1\right)}{2\left(2\cos\left(\frac{\pi}{3}\right) - 1\right)}$
\item $\displaystyle \frac{1}{\arctan 2 + \arctan 3 - \frac{3\pi}{4}} = \frac{1}{\arctan 2 + \arctan 3 - \frac{3\pi}{4}}$
\item $\displaystyle \frac{1}{4\cos^{2}\left(\frac{\pi}{12}\right) - 2 - \sqrt{3}} = \frac{1}{4\cos^{2}\left(\frac{\pi}{12}\right) - 2 - \sqrt{3}}$
\item $\displaystyle \ln\left|2\cos\left(\frac{\pi}{3}\right) - 1\right| = \ln\left|2\cos\left(\frac{\pi}{3}\right) - 1\right|$
\item $\displaystyle \frac{1}{\tan\left(\frac{\pi}{8}\right) + 1 - \sqrt{2}} = \frac{1}{\tan\left(\frac{\pi}{8}\right) + 1 - \sqrt{2}}$
\item $\displaystyle \sum_{k=1}^{n} \frac{1}{\arctan\left(\frac{1}{2}\right) + \arctan\left(\frac{1}{3}\right) - \frac{\pi}{4}} = \frac{n}{\arctan\left(\frac{1}{2}\right) + \arctan\left(\frac{1}{3}\right) - \frac{\pi}{4}}$
\item $\displaystyle \left(\sin\left(\frac{\pi}{6}\right) - \frac{1}{2}\right)^{-n} = \left(\sin\left(\frac{\pi}{6}\right) - \frac{1}{2}\right)^{-n}$
\end{enumerate}
\end{document}
